\documentclass[footheight=20pt, footsepline, headheight=20pt, headsepline]{scrartcl}
%
\usepackage[utf8]{inputenc} % below are various important packages
\usepackage{lmodern}
\usepackage[T1]{fontenc}
\usepackage[english]{babel}
\usepackage{textcomp} 
\usepackage{amsmath}
\usepackage{mathrsfs}
\usepackage{latexsym}
\usepackage{amssymb}	
\usepackage{amsfonts}
\usepackage{theorem}
\usepackage{graphicx}
\usepackage{scrlayer-scrpage}
\usepackage{xcolor}
\usepackage{setspace}
\usepackage{framed}
\usepackage{hyperref} 
\usepackage{pgf,tikz,pgfplots} % possibility to insert geogebra graphs
\usepackage{mathrsfs}
\pgfplotsset{compat=1.15}\usetikzlibrary{arrows} % part of geogebra package
\usepackage{qrcode} % insert qr codes
\usepackage{multicol}
\usepackage{multirow}
\usepackage{xurl}
\usepackage{tabularx}
\usepackage{enumitem}
\usepackage{float}
\usepackage{colortbl,rotating,booktabs} % required packages for Excel2Latex


% Add to length for wider margins
\addtolength{\textwidth}{3cm} % right to margin
\addtolength{\hoffset}{-1.6cm} % left to margin
\addtolength{\voffset}{-2cm} % to top
\addtolength{\textheight}{6.5cm} % to bottom

% Headers-Footers
\definecolor{gro}{gray}{0.6} % define color
\setkomafont{pagehead}{\normalfont\sffamily} % define header
\setkomafont{pagefoot}{\normalfont\sffamily} % define footer
\addtokomafont{headsepline}{\color{gro}} % define header horizontal line
\addtokomafont{footsepline}{\color{gro}} % define footer horizontal line
	\ihead{\color{gro} Athénée de Luxembourg} % header (i=inner=left)
	\chead{\color{gro} Write your title here} % header (c=center)
	\ohead{\color{gro} Your name} % header (o=outer=right)
	\ifoot{\color{gro} } % footer (i=inner=left)
	\cfoot{\color{gro} - {\textbf\thepage} -} % footer (c=center)
	\ofoot{\color{gro} } % footer (o=outer=right)


\renewcommand{\familydefault}{\sfdefault} % font
\linespread{1.2} % increase line spacing

%---------------------------------------------------------------------------
\begin{document} % every document starts with \begin{document}
%\doublespacing
\title{Write your title here \\[1cm] \large{\textbf{Write your research question here}}\\[1cm]} 
\author{Your name}
\date{\vspace{10cm} Athénée de Luxembourg \\[0.5cm] Mathematics - your class \\[0.5cm] Investigation/IA/EE \\[0.5cm] Word Count: xx}
\maketitle
%---------------------------------------------------------------------------
\newpage
\tableofcontents
%---------------------------------------------------------------------------
\newpage
\section{Introduction}
With the purpose of discovering the different fields of study within the umbrella term mathematics, I came across numerical analysis which sparked my interest as it implements algorithms with the aim of solving math- ematical analysis problems with continuous variables. After some further research, I discovered the Weierstrass approximation theorem, defining the key element of polynomial approximation being that \textit{‘any continuous func- tion on a closed and bounded interval can be uniformly approximated on that interval by polynomials to any degree of accuracy.’} \footnote{https://math.stackexchange.com/questions/tagged/weierstrass-approximation} Essentially, this states that a graph which consists of a single unbroken curve in a defined domain can be expressed by a polynomial function whose precision in approximating is determined by its degree. With the aim of analyzing the application of this theorem, I will consider a continuous function which I will approximate with three different polynomial approximation methods.
%---------------------------------------------------------------------------
\section{A few lines of calculations}
The formula of a slope can be proven as follows. For the points $A(x_1;y_1)$ and $B(x_2;y_2)$, that satisfy the equation, I establish:\begin{center}
$\begin{cases}
y_1=mx_1+b\\
y_2=mx_2+b
\end{cases}$
\end{center}
Thus, I can subtract one of these from the other:
\begin{eqnarray*}
&y_2-y_1 &= mx_1+b-(mx_2+b)\\
\Leftrightarrow& y_2-y_1 &= mx_1-mx_2\\
\Leftrightarrow& m &= \dfrac{y_2-y_1}{x_2-x_1}
\end{eqnarray*}
This can be rewritten to find the formula, I used, of a straight line:
\begin{eqnarray*}
y_2&=&y_1+\dfrac{y_2-y_1}{x_2-x_1}\cdot(x_1-x_2)
\end{eqnarray*}
%---------------------------------------------------------------------------
\section{A matrix}
To continue, we transfer our values into an augmented matrix.
$$
\begin{bmatrix}
    7 & 1 & 4 \\
    3 & 1 & 6
\end{bmatrix}
$$
The next step is to add $(-3)$ times the first line to $7$ times the second line to triangulize the matrix.
$$
\Leftrightarrow
\begin{bmatrix}
    21 & 3 & 12\\
    \textcolor{red}{0} & -4 & -30
\end{bmatrix}
$$
%---------------------------------------------------------------------------
\newpage
\section{Inserting a graph from GeoGebra}
Below you can see the graphical representation of $y  = -\dfrac{1}{2} x+ \dfrac{15}{2}$.
\begin{figure}[h!]
\centering
\definecolor{ududff}{rgb}{0.30196078431372547,0.30196078431372547,1}
\resizebox{0.8\linewidth}{!}{%
\begin{tikzpicture}[line cap=round,line join=round,>=triangle 45,x=1cm,y=1cm,scale=1]
\begin{axis}[
x=1cm,y=1cm,
axis lines=middle,
ymajorgrids=true,
xmajorgrids=true,
xmin=-1.5,
xmax=16,
ymin=-1.5,
ymax=9,
xtick={-8,-7,...,21},
ytick={-5,-4,...,13},]
\clip(-8.88,-5.33) rectangle (21.76,13.77);
\draw [line width=2pt,domain=-8.88:21.76] plot(\x,{(--7.5-0.5*\x)/1});
\begin{normalsize}
\draw [fill=ududff] (7,4) circle (2.5pt);
\draw[color=ududff] (7.8,4.42) node {$A(7;4)$};
\draw [fill=ududff] (3,6) circle (2.5pt);
\draw[color=ududff] (3.8,6.42) node {$B(3;6)$};
\draw[color=black] (10.45,0.89) node {$y=-\dfrac{1}{2}x+\dfrac{15}{2}$};
\end{normalsize}
\end{axis}
\end{tikzpicture}}
\caption{Linear model representing $y=-\frac{1}{2} x+ \frac{15}{2}$}
\label{fig:GeoGebra_line}
\end{figure}
As you can see in Figure \ref{fig:GeoGebra_line}, the line passes through the two previously chosen points; $A(7;4)$ and $B(3;6)$, which indicates that our model is correct, that we made no mistakes in our calculations.
%---------------------------------------------------------------------------
\section{Creating a table}
Fill in the following divisibility table.
\begin{table}[h!]
    \centering
    \begin{tabular}{c|c|c|c|c|c|c|c|c|c|c|c|c}
         & $2$ & $3$ & $4$& $5$ & $6$ & $7$ & $9$ & $11$ & $15$ \\ \hline
        $3\,465$    & & & & & & & & & \\ \hline
        $4\,725$    & & & & & & & & & \\ \hline
    \end{tabular}
    \caption{Divisibility table}
    \label{tab:my_table_1}
\end{table}
\begin{center}

\end{center}
%---------------------------------------------------------------------------
\section{Inserting QR Codes}
\begin{multicols}{2}
\begin{center} 
How to install and use the Excel add in \emph{Excel2Latex} \\[.2cm]
    \qrcode{https://express.adobe.com/video/1Dx733VtuxUw9}

YouTube video: Secant Method | Oscar Veliz\\[.2cm]
    \qrcode{https://youtu.be/_MfjXOLUnyw?t=65}
\end{center}
\end{multicols}
%---------------------------------------------------------------------------
\newpage
\section{Inserting an image}
\begin{figure}[h!]
    \centering
    \includegraphics[scale=.07]{samuel-ferrara-unsplash.jpg}
    \caption{Weisshorn, Randa, Switzerland. Author: Samuel Ferrara}
    \label{fig:my_label}
\end{figure}
%---------------------------------------------------------------------------
\section{Customizing the enumerate environment} 
\begin{multicols}{2}
\begin{enumerate}[label=\Alph*.]
    \item Typesetting
    \begin{enumerate}[label=\arabic*)]
        \item Written in \LaTeX~and includes
        \item Cover page
        \item Header/footer
        \item Table of contents
        \item Introduction and conclusion
%        \item No textbook references: Explain what you are doing.\\
    \end{enumerate}
    \item Format
    \begin{enumerate}[label=\alph*)]
        \item Divide into different sections/subsections
        \item Each mathematical expression inside math mode (\$ \$ or \$\$ \$\$)
        \item Each figure should have a caption
        \item Reference figures correctly in your text
%        \item Make sure each label on your figure is readable when printed\\
    \end{enumerate}
\end{enumerate}
\end{multicols}
\section{Customizing the itemize environment}
    \begin{itemize}
        \item First draft: 31-Feb-2029 9pm via MS Teams Assignment
        \item[$\circ$] Final report: 31-Sep-2029 9pm via MS Teams Assignment
        \item[$\ast$] Submit as a PDF file via Teams
        \item[$-$] Name of draft: Name-FirstName-Draft.pdf
        \item[$\square$] Name of final report: Name-FirstName-Final.pdf
    \end{itemize}
%---------------------------------------------------------------------------
\newpage
\section{Bibliography}
\begin{itemize}
    \item Stephenson, M. (2020, December 31). How to calculate darts averages [formulas, chart, pro tips]. DartsGuide. \url{https://dartsguide.net/guides/how-to-calculate-darts-average/}. Accessed 28 May 2021 
    \item Phillips, J. (2018, June 12). Stats analysis: To switch or not to switch? PDC. \url{https://www.pdc.tv/news/2018/06/12/stats-analysis-switch-or-not-switch}. Accessed 28 May 2021 
    \item Chen, L. (2020, July 8). How to interpret and calculate "X Times more likely" statistics. Towards Data Science. \url{https://towardsdatascience.com/how-to-interpret-and-calculate-x-times-more-likely-statistics-daf/538a9e0f4}. Accessed 27 May 2021 
    \item The poisson distribution. \url{https://www.le.ac.uk/users/dsgp1/COURSES/LEISTATS/poisson.pdf}. Accessed 28 November 2021 
\end{itemize}

\end{document}
